
\textbf{Status:} Normative implementation companion to Working Paper v0.3\\
\textbf{Date:} 21 September 2026

\subsection{Repository blueprint}\label{repository-blueprint}

\begin{lstlisting}
where-does-the-function-go/
+--- pyproject.toml
+--- uv.lock / requirements lock
+--- README.md
+--- configs/
|   +--- p1_h_neuron.yaml
|   +--- p2_mask_t0.yaml
|   `--- p3_recoverability.yaml
+--- schemas/
|   +--- run_manifest.schema.json
|   +--- mask.schema.json
|   +--- metric.schema.json
|   `--- trajectory.schema.json
+--- src/wdfg/
|   +--- cli.py
|   +--- config.py
|   +--- provenance.py
|   +--- datasets/
|   |   +--- h_neuron.py
|   |   `--- synthetic_recoverability.py
|   +--- models/
|   |   +--- load.py
|   |   `--- llama_contract.py
|   +--- source_adapters/
|   |   +--- h_neurons.py
|   |   `--- jacobian_lens.py
|   +--- interventions/
|   |   +--- mlp_mask.py
|   |   `--- activation_scale.py
|   +--- capture/
|   |   +--- residual.py
|   |   `--- logits.py
|   +--- stages/
|   |   +--- p1_replication.py
|   |   +--- p2_time_zero.py
|   |   `--- p3_recoverability.py
|   +--- scoring/
|   |   +--- faith_eval.py
|   |   +--- synthetic.py
|   |   `--- gate0.py
|   `--- io/
|       +--- jsonl.py
|       `--- artifacts.py
`--- tests/
    +--- test_llama_contract.py
    +--- test_mask.py
    +--- test_gate0.py
    +--- test_t0_immutable.py
    +--- test_pre_stop.py
    +--- test_continuation_seed.py
    `--- test_synthetic_parser.py
\end{lstlisting}

\subsection{Llama mask contract}\label{llama-mask-contract}

Expected target module per transformer block:

\passthrough{\lstinline!model.model.layers[i].mlp!}

Expected projections:

\begin{itemize}
\tightlist
\item
  \passthrough{\lstinline!gate\_proj!}
\item
  \passthrough{\lstinline!up\_proj!}
\item
  \passthrough{\lstinline!down\_proj!}
\end{itemize}

Expected conceptual forward:

\begin{lstlisting}[language=Python]
u = act_fn(gate_proj(x)) * up_proj(x)
y = down_proj(u)
\end{lstlisting}

Intervention point: input to \passthrough{\lstinline!down\_proj!}.

Preferred mechanism:

\begin{lstlisting}[language=Python]
handle = mlp.down_proj.register_forward_pre_hook(mask_pre_hook)
\end{lstlisting}

Conceptual hook:

\begin{lstlisting}[language=Python]
def mask_pre_hook(module, args):
    (u,) = args
    mask = registry.mask_for_layer(layer_idx, device=u.device, dtype=u.dtype)
    return (u * mask,)
\end{lstlisting}

The actual code must handle the library's exact hook signature and verify shape.

Mask tensor shape should broadcast to:

\passthrough{\lstinline![batch, sequence, intermediate\_size]!}

Recommended stored mask shape:

\passthrough{\lstinline![intermediate\_size]!}

Do not mutate \passthrough{\lstinline!down\_proj.weight!} to impose the persistent constraint.

\subsection{Model contract checks}\label{model-contract-checks}

At startup:

\begin{itemize}
\tightlist
\item
  resolve \passthrough{\lstinline!intermediate\_size!};
\item
  verify each requested neuron index is in range;
\item
  verify target layer exists;
\item
  verify the MLP exposes the expected projections;
\item
  run a one-example identity-mask test;
\item
  abort if contract fails.
\end{itemize}

Pilot runs disable \passthrough{\lstinline!torch.compile!}.

\subsection{Source adapter rule}\label{source-adapter-rule}

P1 H-Neuron identification calls a pinned upstream implementation where possible.

Adapter outputs normalized records:

\begin{lstlisting}
example_id
split
layer_index
neuron_index
cett_score
classifier_weight
is_h_neuron
source_commit
\end{lstlisting}

A local optimized implementation is not allowed to replace the source adapter until equivalence fixtures pass.

\subsection{Run commands}\label{run-commands}

Target CLI surface:

\begin{lstlisting}
wdfg p1 identify --config configs/p1_h_neuron.yaml
wdfg p1 causal   --config configs/p1_h_neuron.yaml --parent-run <id>
wdfg p2 t0       --config configs/p2_mask_t0.yaml --mask <mask.json>
wdfg p3 collect  --config configs/p3_recoverability.yaml
wdfg p3 probe    --config configs/p3_recoverability.yaml --parent-run <id>
wdfg validate-run <run_dir>
\end{lstlisting}

Commands are illustrative API requirements; exact CLI framework is implementation choice.

\subsection{No implicit defaults for scientific parameters}\label{no-implicit-defaults-for-scientific-parameters}

The program may have software defaults for paths.

It must not silently default any scientific parameter that affects the protocol.

Required config keys must fail validation if omitted.

Examples:

\begin{itemize}
\tightlist
\item
  model revision;
\item
  seed;
\item
  coefficient threshold;
\item
  intervention alpha;
\item
  sample IDs/fingerprint;
\item
  K;
\item
  B;
\item
  decoding parameters.
\end{itemize}

\subsection{P1 outputs}\label{p1-outputs}

Required:

\begin{lstlisting}
artifacts/h_neurons.json
artifacts/candidate_scores.parquet
artifacts/faith_eval_predictions.parquet
artifacts/control_predictions.parquet
artifacts/gate0_metrics.json
\end{lstlisting}

\passthrough{\lstinline!h\_neurons.json!} is converted to a canonical mask artifact only after Gate 0 evaluation.

\subsection{P2 outputs}\label{p2-outputs}

Required:

\begin{lstlisting}
artifacts/mask.json
artifacts/mask_audit.json
artifacts/t0_candidate_loci.parquet
artifacts/t0_causal_effects.parquet
\end{lstlisting}

\passthrough{\lstinline!t0\_*!} files are write-once.

The software should refuse overwrite unless an explicit \passthrough{\lstinline!--development-override!} flag is present. Such a run is automatically marked non-result-bearing.

\subsection{P3 natural trajectory representation}\label{p3-natural-trajectory-representation}

Each natural trajectory record contains:

\begin{itemize}
\tightlist
\item
  example ID;
\item
  task family;
\item
  task parameters;
\item
  prompt hash;
\item
  prompt token IDs;
\item
  natural output token IDs;
\item
  parsed natural answer;
\item
  correctness;
\item
  pre-stop prefix token IDs;
\item
  EOS token ID;
\item
  stop-token logit/probability;
\item
  top-1 token ID;
\item
  entropy;
\item
  residual artifact pointer;
\item
  manifest run ID.
\end{itemize}

Large residual tensors are stored outside JSONL.

\subsection{P3 continuation representation}\label{p3-continuation-representation}

Each continuation record contains:

\begin{itemize}
\tightlist
\item
  parent example ID;
\item
  intervention ID (\passthrough{\lstinline!I0!}, \passthrough{\lstinline!I1!}, or \passthrough{\lstinline!I2!});
\item
  k index;
\item
  deterministic sampling seed;
\item
  generated token IDs;
\item
  generated text artifact pointer or text;
\item
  last post-continuation \passthrough{\lstinline!FINAL:!} parse;
\item
  rescue boolean;
\item
  parse status.
\end{itemize}

\subsection{P3 continuation regimes}\label{p3-continuation-regimes}

The implementation must preserve three distinct continuation regimes:

\begin{itemize}
\tightlist
\item \passthrough{\lstinline!I0!}: natural-support continuation from the exact pre-stop prefix, EOS allowed.
\item \passthrough{\lstinline!I1!}: forced-compute continuation, EOS constrained for the configured window.
\item \passthrough{\lstinline!I2!}: procedural re-check with the fixed versioned instruction and no new task facts.
\end{itemize}

The run manifest records the regime explicitly. Analysis code must not pool them into one rescue probability by default.

The continuation count \passthrough{\lstinline!K!} is not hard-coded in v0.6; config validation requires an explicitly frozen value from the post-review precision/cost register.

\subsection{Synthetic generator contract}\label{synthetic-generator-contract}

Generator functions must return:

\begin{lstlisting}[language=Python]
{
  "example_id": "...",
  "family": "integer_arithmetic",
  "parameters": {...},
  "prompt": "... FINAL: <answer> ...",
  "ground_truth": "...",
  "checker_version": "..."
}
\end{lstlisting}

Task difficulty parameters are part of the dataset fingerprint.

Train/val/test generator seeds must not overlap.

\subsection{Dataset fingerprint}\label{dataset-fingerprint}

For each dataset/split, fingerprint:

\begin{itemize}
\tightlist
\item
  canonical list of example IDs;
\item
  normalized prompt text;
\item
  task parameters if generated;
\item
  ground-truth labels;
\item
  preprocessing version.
\end{itemize}

Store SHA-256 of the canonical JSON serialization.

\subsection{Gate 0 evaluator}\label{gate-0-evaluator}

Pure function inputs:

\begin{itemize}
\tightlist
\item
  baseline FaithEval predictions;
\item
  H hard suppression predictions;
\item
  contribution-matched control predictions;
\item
  PPL metrics;
\item
  held-out classifier metrics.
\end{itemize}

Output:

\begin{lstlisting}
{
  "G0_A": "PASS|FAIL",
  "G0_B": "PASS|FAIL",
  "G0_C": "PASS|FAIL",
  "G0_D": "PASS|FAIL",
  "overall": "PASS|FAIL|AMBIGUOUS",
  "evidence": {...}
}
\end{lstlisting}

No manual override in a result-bearing run.

\subsection{Integrity checks for pre-stop recomputation}\label{integrity-checks-for-pre-stop-recomputation}

Natural-run pre-stop record stores:

\begin{itemize}
\tightlist
\item
  pre-stop prefix;
\item
  original top-1 token;
\item
  original EOS probability;
\item
  original logits hash if retained.
\end{itemize}

Recompute from prefix.

Pass if:

\begin{itemize}
\tightlist
\item
  top-1 token is identical;
\item
  absolute EOS-probability difference \textless= configured tolerance.
\end{itemize}

Suggested pilot tolerance: \passthrough{\lstinline!1e-4!} after probability computation in float32.

If fail:

\begin{itemize}
\tightlist
\item
  record \passthrough{\lstinline!FAIL\_DATA\_INTEGRITY!} for that example;
\item
  exclude from recoverability labels;
\item
  do not silently continue.
\end{itemize}

\subsection{Hidden-state capture}\label{hidden-state-capture}

Default:

\begin{itemize}
\tightlist
\item
  last-prefix-token residual state at every transformer layer;
\item
  cast/storage dtype recorded;
\item
  stored using \passthrough{\lstinline!safetensors!} or another deterministic tensor format;
\item
  index file maps example ID -\textgreater{} tensor artifact.
\end{itemize}

Do not serialize arbitrary Python pickle objects for scientific artifacts.

\subsection{J-lens adapter}\label{j-lens-adapter}

J-lens is optional in first P3 pass.

Adapter must not change the underlying P3 trajectory collection.

Feature extraction runs as a downstream analysis against frozen trajectory IDs.

This ensures a failed J-lens integration cannot alter the baseline recoverability dataset.

\subsection{Concurrency}\label{concurrency}

Parallelism must not change random results.

Continuation seed derivation uses example/intervention/k identity, not global RNG order.

All stage outputs are sorted by stable example ID before hashing.

\subsection{Error handling}\label{error-handling}

Scientific failures and software failures are separate.

Examples:

\begin{itemize}
\tightlist
\item
  Gate 0 fail = scientific/protocol result.
\item
  Missing mask coordinate = implementation failure.
\item
  Dataset hash mismatch = data-integrity failure.
\item
  CUDA OOM = hardware failure.
\end{itemize}

Only the first category contributes directly to scientific interpretation.

\subsection{Security / privacy}\label{security-privacy}

Pilot datasets should contain no personal or sensitive user data.

Prompt artifacts should be inspectable and releasable where licensing permits.

\subsection{Acceptance criterion for implementation v0.1}\label{acceptance-criterion-for-implementation-v0.1}

The implementation is ready for P1 result-bearing execution when:

\begin{itemize}
\tightlist
\item
  schema validation passes;
\item
  unit tests pass;
\item
  identity mask matches baseline;
\item
  targeted mask zeros intended coordinates;
\item
  H-Neuron source adapter fixture passes;
\item
  dataset fingerprint is stable across reruns;
\item
  run manifest is complete;
\item
  Gate 0 evaluator passes fixture tests;
\item
  dirty worktree is false.
\end{itemize}
